\begin{center}
    \label{fig:shearstrain}
    % Figure 2: Shear Strain (Simple Shear)
    \begin{tikzpicture}[>=Stealth, thick, scale=1.2]
        % Reference configuration (dashed purple rectangle)
        \draw[red, dashed, line width=1.5pt] (0,0) rectangle (2.5,2);
        
        % Current configuration (parallelogram with horizontal top and bottom)
        \draw[blue, line width=1.5pt] (0,0) -- (2.5,1) -- (2.5, 3) -- (0,2) -- cycle;
        
        % Horizontal dimension - T (from origin to bottom right corner)
        \draw[<->, >=Stealth] (0,-0.5) -- (2.5,-0.5) node[midway, below] {$T$};
        
        % Vertical dimension - D (from x-axis to bottom right corner)
        \draw[<->, >=Stealth] (2.7,0) -- (2.7,1) node[midway, right] {$D$};
        
        % Axes
        \draw[->, very thick] (-0.5,0) -- (4.0,0) node[right] {$x$};
        \draw[->, very thick] (0,-0.3) -- (0,2.8) node[above] {$y$};

        % Label
        \node[below] at (2,-1) {\textbf{Shear Strain}};
    \end{tikzpicture}
    
\end{center}